\onlyinsubfile{\setcounter{section}{5}}
\section{Conclusion}
\label{sec:concl}

\par There is much empirical evidence of a positive relationship between trust and economic performance, at the aggregate and individual level. The results of this paper are in line with this literature. In particular, the RAND longitudinal product version of the HRS contains enough information to create a measure of returns to net wealth. Respondent-level returns exhibit significant heterogeneity, and the persistent component of individual returns is estimated and comparable to the distribution across indivdiuals using the Norwegian population data  \cite{aflgdmlp20}. Moreover, the returns to net wealth exhibit strong evidence of a moderate amount of trust being associated with the maximal level of returns, as \cite{jbpglg2016} shows for household income and trust.

\par The remaining sections of this paper discuss findings relevant to the use of this HRS RAND product that are excluded from the main analysis but can be found along with the rest of the \href{https://github.com/dedwar65/Chp2-TrustandReturns/tree/main/Supplementary}{supplementary material for this project}.

\subsection{Drawbacks and limitations}

\par A first limitation is that the HRS variables used in this paper do not all come from the same source. Most of the demographic, income, and wealth variables come from the cleaned and harmonized \href{https://hrs.isr.umich.edu/news/data-announcements/rand-hrs-longitudinal-file-2018-v1}{RAND HRS Longitudinal File}, but net investment flows must be imported wave-by-wave from the \href{https://hrsdata.isr.umich.edu/data-products/public-survey-data}{HRS core public files}. In the later surveys, these asset-change questions appear in the questionnaire as Section~R (for example, the 2020 module is \href{https://hrsonline.isr.umich.edu/modules/meta/2020/core/qnaire/online/18hr20R.pdf}{Section~R: Asset Change}). This creates a discrepancy in coverage between the standard RAND variables and the flow variables. To partly address this, I use the associated asset-change flags to set some missing flow amounts equal to zero when respondents explicitly report no purchase, sale, contribution, or withdrawal since the previous interview. This increases the usable flow sample, but the flow variables remain much thinner than the core RAND measures.

\par A second limitation is that the trust variables are only observed in 2020. This cannot be overstated. For general trust, there are 900 unique respondents with nonmissing values in the cross section, which is much smaller than the sample available for the standard RAND demographic and wealth variables. In the panel setting, the overlap between nonmissing returns to net wealth and nonmissing general trust is 4,115 person-years in total, ranging from 176 observations in 2002 to 565 in 2020. As a result, the trust analysis is much more sample-constrained than the analogous descriptive and regression exercises for returns, wealth, or demographics.

\par A third limitation concerns the implementation of the correlated random effects (CRE) specification. In this paper, the CRE model is implemented manually by augmenting the random-effects estimator (\texttt{xtreg \ldots, re} in Stata) with respondent-level means of the time-varying regressors---the same Mundlak device that Stata~19 now automates through its \texttt{xtreg, cre} option and the \texttt{estat mundlak} post-estimation command. In earlier Stata workflows, the researcher generated these panel means by hand and tested their joint significance directly; the implementation in this paper follows that convention and is econometrically equivalent in substance. Although survey-year fixed effects enter the model in levels, respondent-level averages of the year indicators are not included in the Mundlak augmentation because the panel is strongly balanced: the within-household means of year dummies are therefore constant across households and collinear with the intercept. The Mundlak means included are consequently limited to time-varying regressors that exhibit meaningful cross-household variation over time, such as wealth deciles, labor-force participation, and marital status.

\subsection{Robustness checks}

\par Labor income does not exhibit the same robust hump-shaped relationship as total income. In the conclusion regressions, the joint test on the trust terms is insignificant for average labor income in both the baseline and extended specifications, and it is only marginal in the 2020 cross section once the full set of controls is added. This weaker pattern is consistent with the HRS sample itself, which oversamples older and retired households and is therefore less naturally suited to labor-income variation than to broader measures of household resources.

\par Within narrower asset classes and portfolio definitions, the trust relationship is not robust in the way it is for returns to net wealth. In the cross section, the uncontrolled specifications for core, IRA, residential, and core-plus-IRA returns sometimes show a hump-shaped pattern, but those results disappear once the baseline controls are added. In that sense, the evidence outside net wealth is best viewed as suggestive at most rather than stable across specifications.

\par Even for returns to net wealth, the other trust measures generally do not reproduce the main result. The individual domain measures, the government-trust aggregates, and the simple average of trust in financial institutions are not jointly significant in the corresponding average-return regressions. The only partial exception is the first principal component for trust in financial institutions, which is weakly significant in the baseline specification, but that evidence disappears once the extended controls are included.

\par Financial literacy adds some information, but it does not overturn the trust results. In the recent extensions, the inflation item matters in the average-return specification, and the literacy score is significant in the FE second-stage regressions for the persistent component. In the CRE model, adding the literacy items one at a time still leaves the trust terms significant for the interest and inflation questions, although broader literacy bundles are less stable. Taken together, financial literacy explains some additional variation in returns, but it does not subsume the main trust relationship.

\par I also explored an instrumental-variables interpretation using depression-based measures and an indicator for non-Hispanic Black as excluded instruments for trust. The first stages are relevant in the sense that these variables predict trust, and the Kleibergen-Paap underidentification tests reject the null that the excluded instruments are irrelevant for trust. However, the endogeneity tests do not reject exogeneity of trust, the weak-identification diagnostics are only borderline, and the weak-instrument-robust Anderson-Rubin tests remain insignificant across the main specifications, meaning the IV exercises do not provide robust evidence that instrumented trust has a causal effect on returns. When the model is overidentified, the overidentification restrictions are not rejected, but the second-stage trust coefficients remain imprecise and unstable. Taken together, these exercises do not provide persuasive evidence for a causal 2SLS interpretation of the effect of trust on returns to net wealth.
